\section*{Chapitre \ref{ch:g1:the-five-senses} --- Les cinq sens}

\begin{solution}{exo:g1:the-five-senses:1}
La vue --- les \omterm{def:g1:the-five-senses:sense}{yeux}~; l'ouïe --- les \omterm{def:g1:the-five-senses:sense}{oreilles}~; l'odorat --- le
nez~; le goût --- la langue~; le toucher --- la \omterm{def:g1:the-five-senses:sense}{peau}.
\end{solution}

\begin{solution}{exo:g1:the-five-senses:2}
(a) Les \omterm{def:g1:the-five-senses:sense}{yeux}. (b) La \omterm{def:g1:the-five-senses:sense}{peau}. (c) Le nez. (d) Les \omterm{def:g1:the-five-senses:sense}{oreilles}.
\end{solution}

\begin{solution}{exo:g1:the-five-senses:3}
La \omterm{def:g1:the-five-senses:sense}{peau}, l'organe du toucher. Elle couvre tout le corps, de la tête
aux pieds.
\end{solution}

\begin{solution}{exo:g1:the-five-senses:4}
La \omterm{def:g1:the-five-senses:sense}{peau} des mains a envoyé le message (froid, dur, dentelé)~; le
cerveau l'a compris et a nommé l'objet.
\end{solution}

\begin{solution}{exo:g1:the-five-senses:5}
Le goût et l'odorat. Le nez fait la moitié du travail de la saveur,
et c'est pourquoi les aliments semblent fades quand il est bouché.
\end{solution}

\begin{solution}{exo:g1:the-five-senses:6}
Par exemple~: ne jamais regarder le Soleil en face --- protège les
\omterm{def:g1:the-five-senses:sense}{yeux}~; garder les sons doux --- protège les \omterm{def:g1:the-five-senses:sense}{oreilles}~; souffler sur
les plats chauds --- protège la langue~; ne rien enfoncer dans
l'\omterm{def:g1:the-five-senses:sense}{oreille} ni le nez --- protège \omterm{def:g1:the-five-senses:sense}{oreilles} et nez~; une bonne lumière
pour lire --- protège les \omterm{def:g1:the-five-senses:sense}{yeux}.
\end{solution}

\begin{solution}{exo:g1:the-five-senses:7}
Parce qu'un son trop fort abîme les \omterm{def:g1:the-five-senses:sense}{oreilles}, et qu'une \omterm{def:g1:the-five-senses:sense}{oreille}
abîmée peut le rester. Les protections gardent le bruit des machines
assez doux pour toute une vie d'écoute.
\end{solution}

\begin{solution}{exo:g1:the-five-senses:8}
Réponse libre --- la plupart des gens en nomment deux ou trois sur
trois. Le nez reconnaît bien les aliments familiers, même sans
l'aide des \omterm{def:g1:the-five-senses:sense}{yeux}.
\end{solution}

\begin{solution}{exo:g1:the-five-senses:9}
L'\omterm{def:g1:the-five-senses:sense}{œil} ne fait que \emph{voir} --- il récolte la lumière et envoie un
message. C'est le cerveau qui comprend~: il reçoit le message et dit
«~c'est mamie~», «~c'est mon école~». Les sens récoltent~; le
cerveau comprend.
\end{solution}

\begin{solution}{exo:g1:the-five-senses:10}
Lire~: du bout des doigts --- le toucher remplace la vue, en sentant
des points en relief sur la page. Traverser la rue~: avec les
\omterm{def:g1:the-five-senses:sense}{oreilles} --- en écoutant les voitures qui arrivent et qui s'en vont
--- souvent avec une canne dont le bout laisse le toucher vérifier
le sol devant.
\end{solution}
